\documentclass[11pt]{article}
\usepackage[a4paper,margin=22mm]{geometry}
\usepackage[T1]{fontenc}
\usepackage{lmodern}
\usepackage{booktabs,tabularx,amsmath}
\usepackage{xcolor}
\definecolor{accent}{HTML}{2459A8}
\usepackage[colorlinks=true,linkcolor=accent,urlcolor=accent]{hyperref}
% Keep saved PDF descriptive properties and provenance fields blank.
\hypersetup{pdftitle={},pdfauthor={},pdfsubject={},pdfkeywords={},pdfcreator={},pdfproducer={}}
\ifdefined\pdfinfoomitdate\pdfinfoomitdate=1\fi
\ifdefined\pdftrailerid\pdftrailerid{}\fi
\ifdefined\pdfsuppressptexinfo\pdfsuppressptexinfo=15\fi
\setlength{\parindent}{0pt}
\setlength{\parskip}{5pt}
\setlength{\tabcolsep}{5pt}
\newcommand{\cond}[1]{\texttt{#1}}
\title{\textbf{Model anonymization through rewriting}\\[4pt]\large Attribution experiment v2: results}
\author{}
\date{October 9, 2026}
\begin{document}
\maketitle
\begin{abstract}
\noindent
Across 160 prompts and five local language models, the strongest tested adaptive
attacker identifies the source model in 92.1\% of unprotected answers. Three
rewrites after persona framing reduce this to 30.5\% (95\% CI 27.5--33.5\%),
while three plain rewrites reach 29.8\% (26.8--32.6\%). Chance is 20\%.
The paired comparison finds no clear difference between the two three-hop
chains, and plain rewrites retain higher embedding similarity and longer
answers. Higher rewrite temperature helps in the one-hop comparison. Neither
a stronger fine-tuned classifier nor a benefit from a fixed rewriter is
established by the specified tests. These results support substantial
reduction in attribution, with residual leakage and unmeasured task correctness.
\end{abstract}

\section{Scope and experimental setup}
This report records the completed v2 run from
\path{experiments/data/results_v2.json}, using the design and hypothesis rules
in \path{experiments/PROTOCOL_v2.md}. It reports saved results without rerunning
collection or model training. The older v1 report is a separate experiment;
\cond{z\_mode} and \cond{para\_1} were not expanded or re-evaluated in v2.

The study contains 800 units: 160 prompts times five choices of source model
$Z$. The original 80 prompts are supplemented with 80 prompts across the same
seven task types. Each model answers every prompt. The candidate models are
Meta's Llama 3.2 3B, Google's Gemma 3 4B, Alibaba's Qwen3 4B 2507,
Microsoft's Phi-4-mini, and Mistral's Ministral 3 3B, run locally in LM Studio.

A circuit draws three distinct relay labs uniformly without replacement; $Z$'s
lab can also be a relay. Rewrites run in exit, middle, guard order. Paired
\cond{nodes\_k} and \cond{plain\_k} conditions use the same circuit for a unit.
Sampling uses temperature 0.7 and a 400-token output limit, except the hot
one-hop rewrite, which uses temperature 1.2. All evaluated final texts pass
through the same lab-name and self-identification scrubber.

The completed cache contains 8,880 unique calls, including reused v1 calls and
v1-only conditions. None of the cached outputs is empty. The v2 analysis log
ends with a successful results write after 26,824 seconds (about 7.45 hours).

\clearpage
\section{Conditions and attackers}
\begin{center}
\small
\begin{tabularx}{\linewidth}{lX}
\toprule
Condition & Text observed by the attacker \\
\midrule
\cond{direct} & $Z$ answers the plain prompt. \\
\cond{persona\_fwd} & $Z$ answers the nested guard/middle/exit persona prompt. \\
\cond{nodes\_1..3} & Persona answer, then one, two, or three independent model rewrites. \\
\cond{plain\_1..3} & Plain answer, then the same one, two, or three relay models. \\
\cond{fixed\_1} & Plain answer rewritten once by Meta for every unit. \\
\cond{plain\_1\_hot} & Same exit relay as \cond{plain\_1}, at temperature 1.2. \\
\bottomrule
\end{tabularx}
\end{center}

\textbf{A1} uses TF-IDF character 2--5-grams and word 1--2-grams with logistic
regression. \textbf{A2} uses nomic-embed-text-v1.5 embeddings with logistic
regression. Both select regularization through internal cross-validation.
\textbf{A3} fine-tunes ModernBERT-base with a five-class head for four epochs,
learning rate $3\times10^{-5}$, weight decay 0.01, batch size 16, maximum
length 256 tokens, 10\% linear warmup, and seed 0. It uses no test-fold tuning
or early stopping.

Adaptive attackers train and test within the same condition. Static attackers
train only on \cond{direct}. Outer evaluation uses five folds grouped by
prompt: each fold trains on 128 prompts (640 texts) and tests on 32 prompts
(160 texts), with no prompt shared across outer train and test sets.
Accuracy CIs use 2,000 bootstrap resamples over prompts. Condition differences
use paired prompt bootstraps.

The best attacker for a condition is the one with the highest observed
adaptive accuracy. Paired condition tests use a single attacker on both sides,
selected by the highest mean adaptive accuracy over the pair. These selections
are part of the stated analysis rule; intervals are not adjusted for attacker
selection or multiple hypothesis testing.

\section{Adaptive attribution accuracy}
Values are percentages, with 95\% CIs in brackets. Lower accuracy means less
successful source identification. Chance is 20\%. Bold identifies the
strongest tested adaptive attacker for each condition.
\begin{center}
\small
\begin{tabular}{lrrr}
\toprule
Condition & A1: TF-IDF & A2: embeddings & A3: ModernBERT \\
\midrule
\texttt{direct} & \textbf{92.1 [90.0, 94.1]} & 61.1 [57.9, 64.4] & 88.4 [85.8, 90.8] \\
\texttt{persona\_fwd} & 84.6 [81.9, 87.2] & 59.8 [56.4, 63.2] & \textbf{85.5 [82.9, 88.1]} \\
\texttt{nodes\_1} & 39.0 [35.8, 42.2] & 34.5 [31.8, 37.2] & \textbf{40.4 [37.2, 43.2]} \\
\texttt{nodes\_2} & 35.0 [32.5, 37.5] & 29.5 [26.9, 32.0] & \textbf{35.6 [32.9, 38.4]} \\
\texttt{nodes\_3} & \textbf{30.5 [27.5, 33.5]} & 27.1 [24.8, 29.5] & 30.2 [27.5, 32.8] \\
\texttt{plain\_1} & 49.0 [45.9, 52.3] & 33.5 [30.9, 36.2] & \textbf{50.2 [46.9, 53.6]} \\
\texttt{plain\_2} & 35.1 [32.1, 38.4] & 27.1 [24.8, 29.6] & \textbf{37.6 [34.4, 40.9]} \\
\texttt{plain\_3} & 25.2 [22.4, 28.1] & 24.8 [22.4, 27.3] & \textbf{29.8 [26.8, 32.6]} \\
\texttt{fixed\_1} & \textbf{48.9 [45.8, 52.4]} & 33.9 [31.2, 36.5] & 43.2 [39.6, 46.9] \\
\texttt{plain\_1\_hot} & 40.4 [37.2, 43.6] & 30.8 [28.1, 33.5] & \textbf{43.5 [40.2, 46.9]} \\
\bottomrule
\end{tabular}
\end{center}

The largest observed reduction occurs when answers are rewritten by relay
models. The best adaptive accuracy falls from 92.1\% on \cond{direct} to
40.4\% on \cond{nodes\_1}, 35.6\% on \cond{nodes\_2}, and 30.5\% on
\cond{nodes\_3}. Plain rewriting also improves with hop count: 50.2\%,
37.6\%, and 29.8\%. These summaries select the best attacker separately for
each condition; the paired tests below hold the attacker fixed.

\clearpage
\section{Hypothesis tests and interpretation}
Differences below are percentage points, with paired 95\% CIs. H13 instead
reports an accuracy and its CI. ``Supported'' follows the protocol's stated
rule; ``not supported'' does not establish that an effect is absent.
\begin{center}
\small
\begin{tabular}{llrl}
\toprule
Test & Comparison & Estimate [95\% CI] & Verdict \\
\midrule
H7 & A3 minus A1 on nodes 3 & -0.25 [-3.38, +3.00] & Not supported \\
H8 & A1: full minus half training data on nodes 3 & +1.87 [-1.00, +4.63] & Not supported \\
H9 (1) & Nodes 1 minus plain 1, A3 & -9.87 [-13.87, -6.00] & Not supported \\
H9 (2) & Nodes 2 minus plain 2, A3 & -2.00 [-5.63, +1.75] & Supported \\
H9 (3) & Nodes 3 minus plain 3, A3 & +0.50 [-3.50, +4.50] & Supported \\
H10 & Plain 3 minus plain 1, A3 & -20.50 [-24.50, -16.62] & Supported \\
H11 & Fixed 1 minus plain 1, A1 & -0.12 [-4.25, +3.75] & Not supported \\
H12 & Hot plain 1 minus plain 1, A3 & -6.75 [-10.50, -3.00] & Supported \\
H13 & Nodes 3 accuracy above chance, A1 & 30.50 [27.62, 33.50] & Supported \\
\bottomrule
\end{tabular}
\end{center}

\textbf{Persona framing.} With A3 held fixed, the persona chain has 9.88 points
lower attribution after one hop than the plain chain (CI $-13.87$ to $-6.00$).
Thus H9's prediction of no detectable difference fails at one hop. The two-
and three-hop differences include zero. H9 is recorded as supported there
under the protocol's rule, but this is not an equivalence test: the three-hop
CI still permits effects from $-3.50$ to $+4.50$ points.

\textbf{More hops.} Three plain rewrites lower A3 accuracy by 20.50 points
relative to one rewrite (CI $-24.50$ to $-16.63$), supporting H10. The three-hop
plain and persona chains have similar observed performance against A3
(29.75\% and 30.25\%).

\textbf{Fixed rewriter and temperature.} The paired fixed-versus-random test
uses A1: the difference is $-0.12$ points (CI $-4.25$ to $+3.75$), so H11 is
not supported. Using temperature 1.2 lowers A3 one-hop accuracy by 6.75 points
(CI $-10.50$ to $-3.00$), supporting H12. The fixed-rewriter conclusion applies
to the tested Meta model and this data.

\textbf{Attacker strength and training size.} A3 does not clearly outperform
A1 on \cond{nodes\_3}: the difference is $-0.25$ points (CI $-3.38$ to
$+3.00$). Doubling A1's training data from half to all available prompts yields
$+1.87$ points (CI $-1.00$ to $+4.63$). H7 and H8 are not supported. The latter
does not prove saturation or rule out improvement with substantially more data.

\textbf{Residual leakage.} H13 is supported: the separately resampled CI for
A1 on \cond{nodes\_3} is 27.63--33.50\%, wholly above 20\% chance. This CI
and the main table's 27.5--33.5\% interval use separate bootstrap draws for
the same 30.5\% estimate.

\clearpage
\section{Similarity, length, and remaining signal}
Utility is the mean cosine similarity of the final text to $Z$'s plain answer,
using nomic embeddings. It is a proxy for answer similarity, not a measure of
correctness or complete fact preservation. Word counts are means rounded to
the nearest word.
\begin{center}
\small
\begin{tabular}{lrr}
\toprule
Condition & Mean cosine & Mean words \\
\midrule
\texttt{direct} & 1.000 & 247 \\
\texttt{persona\_fwd} & 0.891 & 165 \\
\texttt{nodes\_1} & 0.880 & 139 \\
\texttt{nodes\_2} & 0.874 & 126 \\
\texttt{nodes\_3} & 0.869 & 122 \\
\texttt{plain\_1} & 0.952 & 199 \\
\texttt{plain\_2} & 0.933 & 174 \\
\texttt{plain\_3} & 0.918 & 160 \\
\texttt{fixed\_1} & 0.953 & 174 \\
\texttt{plain\_1\_hot} & 0.924 & 212 \\
\bottomrule
\end{tabular}
\end{center}

The saved reference similarity between two different models answering the
same prompt is 0.917; the saved same-model resampling reference is 0.954.
These are descriptive reference values, not hard bounds on utility.

At three hops, plain rewriting has higher mean cosine (0.918 versus 0.869)
and longer answers (160 versus 122 words) than the persona chain. Both shorten
answers relative to the 247-word direct baseline. No confidence intervals or
task-level tests were saved for these utility comparisons, so the observed
similarity difference should not be described as a tested correctness benefit.

Leakage is uneven across source models. With A1 on \cond{nodes\_3}, Mistral is
identified in 87 of 160 cases (54.4\%); Meta, Google, Alibaba, and Microsoft
are identified at 28.8\%, 21.3\%, 21.9\%, and 26.2\%, respectively. For
\cond{plain\_3}, A3's corresponding rates are 29.4\%, 38.8\%, 15.6\%,
35.6\%, and 29.4\%. These are descriptive row accuracies without per-model CIs.
Different strongest attackers are used for the two matrices, so this breakdown
is not a paired test of persona effects for individual models.

\section{Practical conclusions and limits}
Three model rewrites greatly reduce source attribution in this local-model
experiment. Plain three-hop rewriting is a useful simpler design: its tested
adaptive performance is similar to the persona chain, with higher observed
answer similarity and length. Persona framing helps at one hop but has no
clearly detected advantage after two or three hops.

The results do not establish anonymity: three-hop attribution remains above
chance. The study covers five small local models, a finite prompt set, and
three attacker families. Frontier models, stronger attacks, larger training
sets, refusals, and content-based cues may behave differently. Provider
visibility, traffic analysis, timing, watermark robustness, and cryptographic
privacy guarantees are outside the evaluated threat model. The remaining
signal's cause cannot be determined from these classifiers alone.

\section{Source artifacts and verification}
All numeric tables are transcribed from \path{experiments/data/results_v2.json}.
The protocol is \path{experiments/PROTOCOL_v2.md}; implementation is
\path{experiments/analyze_v2.py}; predictions and raw calls are in
\path{experiments/data/preds_v2.json} and \path{experiments/data/calls.jsonl}.
The completed run is recorded in \path{experiments/data/analyze_v2.log}.

As a separate software check, \path{tests/test_blackbox.py} was rerun with
\texttt{python -m pytest -q -p no:cacheprovider}: 11 tests passed. These mock
backend tests verify program behavior; they do not validate the real-model
anonymity or answer correctness claims.

\clearpage
\appendix
\section{Static attackers}
Each attacker is trained on direct answers only. Values are percentages with
95\% prompt-bootstrap CIs. The direct baseline appears in the adaptive table
because its training and test condition are already the same.
\begin{center}
\small
\begin{tabular}{lrrr}
\toprule
Condition & A1: TF-IDF & A2: embeddings & A3: ModernBERT \\
\midrule
\texttt{persona\_fwd} & 63.7 [60.9, 66.6] & 34.8 [32.1, 37.5] & 52.5 [49.5, 55.6] \\
\texttt{nodes\_1} & 33.0 [30.4, 35.6] & 26.6 [24.4, 29.0] & 28.9 [26.5, 31.3] \\
\texttt{nodes\_2} & 28.4 [25.8, 30.9] & 24.4 [22.5, 26.4] & 25.6 [23.6, 27.6] \\
\texttt{nodes\_3} & 25.8 [23.1, 28.2] & 22.2 [20.6, 23.9] & 24.4 [22.2, 26.5] \\
\texttt{plain\_1} & 42.8 [39.4, 46.0] & 33.2 [30.6, 36.0] & 43.8 [40.5, 47.1] \\
\texttt{plain\_2} & 33.9 [30.8, 36.9] & 28.6 [26.1, 31.1] & 32.9 [30.1, 35.6] \\
\texttt{plain\_3} & 28.2 [25.4, 31.2] & 25.2 [23.1, 27.5] & 27.5 [25.0, 30.1] \\
\texttt{fixed\_1} & 41.5 [38.9, 44.1] & 32.8 [30.4, 35.2] & 36.8 [34.2, 39.4] \\
\texttt{plain\_1\_hot} & 37.9 [34.5, 41.1] & 31.4 [28.7, 34.0] & 39.9 [37.2, 42.6] \\
\bottomrule
\end{tabular}
\end{center}

\section{Predictions naming the last relay}
The rate is the percentage of predictions naming the final rewrite lab,
restricted to units where that lab differs from $Z$. It describes recognition
of the relay's voice, not correct recognition of the source. Percentages below
are rounded to one decimal; no CIs for these rates were saved.
\begin{center}
\small
\begin{tabular}{lrrrrrr}
\toprule
Condition & A1 adapt. & A2 adapt. & A3 adapt. & A1 static & A2 static & A3 static \\
\midrule
\texttt{nodes\_1} & 27.4 & 20.4 & 28.3 & 41.3 & 24.5 & 37.0 \\
\texttt{nodes\_2} & 24.7 & 18.2 & 25.8 & 35.0 & 24.1 & 30.8 \\
\texttt{nodes\_3} & 22.4 & 20.2 & 22.8 & 36.2 & 23.0 & 29.6 \\
\texttt{plain\_1} & 22.9 & 21.5 & 24.0 & 47.1 & 26.9 & 35.3 \\
\texttt{plain\_2} & 24.4 & 20.5 & 19.6 & 42.1 & 24.4 & 30.5 \\
\texttt{plain\_3} & 26.0 & 19.7 & 21.8 & 41.0 & 26.3 & 32.6 \\
\texttt{fixed\_1} & 12.0 & 17.5 & 10.3 & 60.6 & 29.4 & 50.9 \\
\texttt{plain\_1\_hot} & 20.9 & 19.0 & 18.6 & 51.7 & 32.0 & 35.1 \\
\bottomrule
\end{tabular}
\end{center}

For example, on \cond{nodes\_3}, A1's static last-relay prediction rate is
36.2\%, compared with 22.4\% for adaptive A1. The fixed Meta rewrite yields
60.6\% for static A1, while its adaptive rate is 12.0\%. The exclusion of
units where the last lab is $Z$ is essential to interpreting these rates.

\clearpage
\section{Training-data learning curves}
Fractions refer to the 128 training prompts in each outer fold. The 25\%,
50\%, and 100\% settings use 32, 64, and 128 prompts (160, 320, and 640
training texts), respectively. All settings test on the same held-out prompts.
Values are percentages with 95\% CIs; the full-data curve intervals can differ
slightly from the main table because of separate bootstrap draws.
\begin{center}
\small
\begin{tabular}{llrrr}
\toprule
Condition & Attacker & 25\% training & 50\% training & 100\% training \\
\midrule
\texttt{direct} & A1 & 82.1 [79.0, 85.1] & 88.6 [86.2, 90.9] & 92.1 [89.9, 94.2] \\
\texttt{nodes\_3} & A1 & 28.4 [25.8, 30.9] & 28.6 [26.1, 31.5] & 30.5 [27.5, 33.5] \\
\texttt{plain\_3} & A1 & 23.9 [21.2, 26.4] & 25.0 [22.2, 27.9] & 25.2 [22.5, 28.0] \\
\texttt{nodes\_3} & A3 & 26.1 [23.8, 28.5] & 29.2 [26.6, 31.9] & 30.2 [27.8, 32.9] \\
\bottomrule
\end{tabular}
\end{center}

The direct A1 baseline rises from 82.1\% to 92.1\% with additional training
prompts. On the three-hop persona chain, the corresponding A1 curve is
28.4\%, 28.6\%, and 30.5\%; A3 is 26.1\%, 29.2\%, and 30.2\%.
Only the specified H8 full-versus-half comparison is reported as a hypothesis
test; these other curve differences are descriptive.

\section{Saved confusion matrices}
Rows are the true source model and columns the predicted model. Each row sums
to 160 cases. The matrices use the best adaptive attacker for their condition.
M = Meta, G = Google, A = Alibaba, P = Microsoft/Phi, R = Mistral.
\subsection*{\texttt{direct} (A1)}
\begin{center}
\small
\begin{tabular}{lrrrrr}
\toprule
True source & M & G & A & P & R \\
\midrule
Meta & 155 & 2 & 1 & 0 & 2 \\
Google & 4 & 146 & 4 & 3 & 3 \\
Alibaba & 9 & 3 & 131 & 0 & 17 \\
Microsoft & 1 & 0 & 2 & 155 & 2 \\
Mistral & 1 & 1 & 7 & 1 & 150 \\
\bottomrule
\end{tabular}
\end{center}
\subsection*{\texttt{nodes\_3} (A1)}
\begin{center}
\small
\begin{tabular}{lrrrrr}
\toprule
True source & M & G & A & P & R \\
\midrule
Meta & 46 & 27 & 29 & 35 & 23 \\
Google & 27 & 34 & 29 & 29 & 41 \\
Alibaba & 38 & 31 & 35 & 30 & 26 \\
Microsoft & 33 & 24 & 32 & 42 & 29 \\
Mistral & 23 & 15 & 18 & 17 & 87 \\
\bottomrule
\end{tabular}
\end{center}
\clearpage
\section*{Confusion matrices (continued)}
\subsection*{\texttt{plain\_3} (A3)}
\begin{center}
\small
\begin{tabular}{lrrrrr}
\toprule
True source & M & G & A & P & R \\
\midrule
Meta & 47 & 27 & 17 & 38 & 31 \\
Google & 16 & 62 & 12 & 25 & 45 \\
Alibaba & 29 & 31 & 25 & 40 & 35 \\
Microsoft & 30 & 21 & 20 & 57 & 32 \\
Mistral & 28 & 39 & 18 & 28 & 47 \\
\bottomrule
\end{tabular}
\end{center}
\subsection*{\texttt{fixed\_1} (A1)}
\begin{center}
\small
\begin{tabular}{lrrrrr}
\toprule
True source & M & G & A & P & R \\
\midrule
Meta & 67 & 30 & 24 & 18 & 21 \\
Google & 13 & 96 & 18 & 15 & 18 \\
Alibaba & 30 & 26 & 47 & 18 & 39 \\
Microsoft & 8 & 20 & 12 & 108 & 12 \\
Mistral & 26 & 24 & 26 & 11 & 73 \\
\bottomrule
\end{tabular}
\end{center}
\end{document}
